\chapter{System Design}


This chapter describes the testbed used to evaluate anti-spam tools under one-way delay representing a deep-space DTN link. It sets out the node layout, the host model and the reasons for it, the tools under test, the DTN configuration, the delay simulation, and the experimental design.

\section{System Architecture}
\label{sec:arch}

The testbed has two networks joined by a relay. Each node runs an ION instance (Section~\ref{sec:background-ion}) and is a separate virtual machine running Ubuntu 22.04 LTS. Earth uses 192.168.100.0/24 and Moon uses 192.168.200.0/24. The nodes are:

\begin{itemize}
  \item \textbf{Mail1 (Earth).} The Earth mail server. Runs Postfix and Dovecot. Address 192.168.100.10.
  ION node 1.
  \item \textbf{Space (Relay).} The transit node between the two networks. Runs an
  ION node and carries the delay and packet-capture configuration. Earth-side
  interface 192.168.100.1, Moon-side interface 192.168.200.1. ION node 2.
  \item \textbf{Mail2 (Moon).} The Moon mail server. Runs Postfix, Dovecot, and
  the three anti-spam tools. Address 192.168.200.11. ION node 3.
  \item \textbf{Client1 (Earth).} User terminal on the Earth network. Runs Thunderbird and a script (the test harness)(Section~\ref{sec:harness}) that injects the spam messages. Connects to Mail1 over IMAP (port 143) and SMTP submission (port 587, STARTTLS).
  \item \textbf{Client2 (Moon).} User terminal on the Moon network. Runs
  Thunderbird and connects to Mail2 the same way.
\end{itemize}

Mail moves server to server across the DTN link: a message submitted on one mail server is carried as a bundle through the relay to the other. ION runs on all three servers, so the bundle path is node 1 (Earth) to node 2 (relay) to node 3 (Moon). The relay is the only path between the two networks, so all traffic between Earth and Moon passes through it. This is where the one-way link delay is applied, standing in for the Earth–Moon propagation time. Mailboxes use mbox format, and each user reads mail from their local server. The topology is shown in Figure~\ref{fig:topology}.

\begin{figure}[ht]
\centering
\resizebox{\textwidth}{!}{%
\begin{tikzpicture}[
  font=\small,
  box/.style={draw, rounded corners, align=center, minimum width=3.2cm, minimum height=1.5cm},
  cl/.style={draw, rounded corners, align=center, minimum width=3.2cm, minimum height=1cm},
  >={Stealth[]}
]
\node[box] (mail1) {\textbf{Mail1 -- Earth}\\Mail Server: Postfix, Dovecot\\IP: 192.168.100.10\\ION node: 1};
\node[cl, below=1.2cm of mail1] (client1) {\textbf{Client1 -- Earth}\\Thunderbird};
\node[box, right=2.6cm of mail1] (relay) {\textbf{Space -- Relay}\\Delay + capture\\IP: 192.168.100.1 / 192.168.200.1\\ION node: 2};

\node[box, right=2.6cm of relay] (mail2) {\textbf{Mail2 -- Moon}\\Mail Server: Postfix, Dovecot\\IP: 192.168.200.11\\ION node: 3};
\node[cl, below=1.2cm of mail2] (client2) {\textbf{Client2 -- Moon}\\Thunderbird};
\draw[<->, thick] (client1) -- node[right, font=\footnotesize]{IMAP / SMTP} (mail1);
\draw[<->, thick] (mail1) -- (relay);
\draw[<->, thick] (relay) -- node[above, font=\footnotesize, align=center]{DTN link\\(delay)} (mail2);
\draw[<->, thick] (client2) -- node[right, font=\footnotesize]{IMAP / SMTP} (mail2);
\end{tikzpicture}%
}
\caption{Testbed topology. Each client connects only to its local mail server. The
relay is the only path between the Earth and Moon networks and is where the one-way
link delay is applied.}
\label{fig:topology}
\end{figure}


\paragraph{Data flow (Earth to Moon).}
The path of a message is shown in Figure~\ref{fig:dataflow}.

\begin{figure}[H]
\centering
\begin{tikzpicture}[
  font=\small, >={Stealth[]},
  flowstep/.style={draw, rounded corners, align=center, minimum width=8.5cm, minimum height=0.85cm},
  node distance=0.45cm
]
\node[flowstep] (s1) {User writes the message in Thunderbird on Client1 (Earth)};
\node[flowstep, below=of s1] (s2) {Client1 submits it to Postfix on Mail1 over SMTP};
\node[flowstep, below=of s2] (s3) {Mail1 hands the message to its ION node (node 1)};
\node[flowstep, below=of s3] (s4) {ION encapsulates the message into a bundle};
\node[flowstep, dashed, below=of s4] (s5) {Bundle is forwarded through Space (node 2) : one-way delay applied};
\node[flowstep, below=of s5] (s6) {ION on Mail2 (node 3) receives the bundle and rebuilds the message (Moon)};
\node[flowstep, below=of s6] (s7) {Postfix and Dovecot on Mail2 deliver it to the recipient's mailbox.};
\node[flowstep, below=of s7] (s8) {User reads it in Thunderbird on Client2 over IMAP};
\foreach \a/\b in {s1/s2,s2/s3,s3/s4,s4/s5,s5/s6,s6/s7,s7/s8} \draw[->] (\a) -- (\b);
\end{tikzpicture}
\caption{End-to-end path of an email from Earth to Moon across the DTN link.}
\label{fig:dataflow}
\end{figure}

\section{Environment and Technology Selection}
\label{sec:env}

\subsection{Virtual machines over containers}
The testbed runs on full virtual machines (Ubuntu 22.04 LTS, ARM64, under UTM on macOS). Containers were considered but not used. They were not used for three reasons:

\begin{itemize}
   \item \textbf{VMs match the topology.} The testbed mimics three separate machines — Earth, Space, Moon — connected by a long-delay link, and full VMs naturally model this, whereas containers share one kernel on one host.
  \item \textbf{Simpler delay setup.} Setting up an intermediate functional block (ifb), the virtual device used to apply delay to incoming traffic (Section~\ref{sec:delay}), is fiddly and error-prone inside containers, but it is straightforward on a VM where each node has its own kernel and network stack.
  \item \textbf{DTN state stays isolated per node.} Each ION node keeps its own bundle store and contact schedule, and a full VM ensures the link delay sits on the path itself rather than leaking into a node's own processing.
\end{itemize}

\section{Email stack}
\begin{itemize}
  \item \textbf{Postfix} was chosen for its clear interface for handing mail to an external program, needed to pass messages to the DTN layer.
  \item \textbf{Dovecot} was chosen because it works directly with mbox format and the local system accounts that Postfix delivers to.
  \item \textbf{Thunderbird}was chosen as a standard mail client, so the message traffic in the testbed matches what a normal user would generate.
\end{itemize}

\section{DTN stack}
\begin{itemize}
  \item \textbf{ION-DTN}was chosen for its deep-space flight heritage and its existing tool for carrying email over the Bundle Protocol, which connects directly to the mail stack.
\end{itemize}

\section{Anti-Spam Tool Selection}
\label{sec:tools}

Three tools were selected, each with a different external dependency:

\begin{itemize}
  \item \textbf{SpamAssassin 3.4.6} --- content scoring that depends on DNS
  blocklist lookups.
  \item \textbf{rspamd 2.7} --- content scoring with a large number of
  asynchronous DNS lookups.
  \item \textbf{Vipul's Razor 2.84} --- collaborative filtering that checks each message against a central catalogue server, which holds signatures of known spam reported by other users.
\end{itemize}

The three cover the main ways a filter reaches the network: blocklist DNS (SpamAssassin, rspamd) and a TCP request to a central reputation server (Razor). Testing one tool would show only one failure path.


\section{DTN Design}
\label{sec:dtn}
The DTN layer uses ION-DTN 4.1.4~\cite{ionsoftware} and BPMail. The three servers run ION as nodes on a single ipn scheme (Section~\ref{sec:background-dtn}): Mail1 (Earth) is node 1, Space (relay) is node 2, and Mail2 (Moon) is node 3. Bundles are carried over a UDP convergence layer (Section~\ref{sec:background-dtn}), with each node listening on its own port: 1113 on Mail1, 2113 on Space, and 3113 on Mail2. The relay forwards bundles between the two mail nodes, so the path from Earth to Moon is node 1 to node 2 to node 3. The contact plan (Section~\ref{sec:background-ion}) sets the windows during which the nodes can exchange bundles. BPMail sits above ION on each mail server: it takes a message from Postfix, encapsulates it into a bundle for the destination node, and rebuilds it on arrival. The full configuration, including the SDR settings (Section~\ref{sec:background-ion}), the contact plan, and the startup order, is given in the implementation chapter.

\section{Delay Simulation Design}
\label{sec:delay}
Delay is applied with netem on the relay, to all traffic crossing the boundary between the Earth and Moon networks. An Earth–Moon link delays travel in both directions: a signal takes about 1.3 seconds each way, so a round-trip takes about 2.6 seconds. The testbed emulates this by delaying each leg on the relay's Moon-facing interface.

The complication is that \texttt{tc} can only shape traffic \emph{leaving} an interface, so the two legs have to be handled separately. The Earth-to-Moon leg is the egress direction on that interface and is delayed by a netem qdisc attached to the interface root. The Moon-to-Earth leg arrives as ingress, and ingress has no egress queue to attach a qdisc to. To delay it, incoming traffic is redirected to an Intermediate Functional Block (ifb) device, whose own egress path can be shaped normally, and the same netem delay is applied there. The two mechanisms, the root qdisc and the ifb device are therefore the two halves of a single symmetric delay, not a way to delay one direction while sparing the other. Traffic that stays within a single network never reaches the relay and is unaffected. This is shown in Figure~\ref{fig:ifb}.



\begin{figure}[ht]
\centering
\resizebox{0.95\textwidth}{!}{%
\begin{tikzpicture}[
  font=\small,
  box/.style={draw, rounded corners, align=center, minimum width=2.8cm, minimum height=1.3cm}
]
\node[box] (earth) {Earth\\node};
\node[box, right=4.2cm of earth] (relay) {Relay (Space)\\netem delay + ifb0};
\node[box, right=4.2cm of relay] (moon) {Moon\\node};

% Earth-to-Moon leg (top row, left to right)
\draw[-latex, very thick] ([yshift=12pt]earth.east) -- node[above, font=\footnotesize, align=center]{Earth-to-Moon leg:\\delayed 1300\,ms (interface egress)} ([yshift=12pt]relay.west);
\draw[-latex, very thick] ([yshift=12pt]relay.east) -- ([yshift=12pt]moon.west);

% Moon-to-Earth leg (bottom row, right to left)
\draw[-latex, very thick] ([yshift=-12pt]moon.west) -- node[below, font=\footnotesize, align=center]{Moon-to-Earth leg:\\delayed 1300\,ms (ifb ingress)} ([yshift=-12pt]relay.east);
\draw[-latex, very thick] ([yshift=-12pt]relay.west) -- ([yshift=-12pt]earth.east);
\end{tikzpicture}%
}
\caption{Symmetric delay. Each leg crossing the relay is delayed by 1300\,ms: the
Earth-to-Moon leg on the interface's egress path, and the Moon-to-Earth leg on an
ifb device, since ingress cannot be shaped directly. The two legs together give a
round-trip time of about 2.6\,s. Traffic that stays within one network does not cross
the relay and is not delayed.}
\label{fig:ifb}
\end{figure}

Delay is the independent variable. The experiments use a representative one-way delay of 1.3 seconds, a round-trip of about 2.6 seconds, as a deep-space operating point. The testbed allows any value to be set, so the same setup covers the range from short Earth–Moon delays to multi-minute deep-space delays.

\subsection{DNS Resolver}

The blocklist rules in every tool under test resolve DNSBL and URIBL names, so
the Moon node needs a working recursive resolver. A large public resolver cannot
serve this role: the major blocklist operators refuse or rate-limit queries that
arrive through open public resolvers such as those run by Google or Cloudflare,
returning an error or an empty answer rather than a listing~\cite{spamhaus}.
Scanning through one would report no spam for reasons unrelated to delay and make
the results meaningless.

The testbed therefore uses a local resolver, running on an off-testbed virtual machine with ordinary internet access to the live blocklist zones. The Moon node runs a local unbound instance that forwards to it, for the following reasons:

\begin{itemize}
  \item \textbf{Valid blocklist answers.} A dedicated resolver querying the blocklist zones is what those operators expect, so listings are returned normally and all that is added to a lookup is the injected link delay.
  \item \textbf{Control of the resolver's own timers.} A public resolver cannot
  be reconfigured, yet the timers that decide how a resolver treats a late answer
  turn out to determine whether detection survives delay (Chapter~5). A private
  resolver can be tuned and its configuration recorded.
  \item \textbf{Reproducibility.} Its cache and per-server state can be reset
  before each run and dumped afterwards, so every run starts cold and its DNS
  state is on record.
\end{itemize}

The recursive resolver sits off the Moon network, so every query the scanner makes crosses the relay and is delayed like the mail.


\section{Summary}
\label{sec:design-summary}

This chapter set out a five-node testbed for measuring anti-spam tool behaviour
under link delay. The design uses full virtual machines for kernel and DNS
isolation, a private resolver to give the Moon side a working DNS path, three tools
chosen to cover the main external dependencies, and an ifb-based scheme to apply
delay. The next chapter describes how this design was built and configured, and the chapter after it presents the results.